\begin{proposition}[Proved under stated hypotheses: local purity along the boundary]
\label{D:local-purity-boundary}
Let $Z$ be a normal connected proper scheme over an algebraically
closed field, let $Y\subsetneq Z$ be closed, and put $U=Z\setminus Y$.
Suppose that for every $y\in Y$ at least one of the following holds:
\begin{enumerate}
\item $\mathcal O_{Z,y}$ is regular of dimension at least two;
\item $\mathcal O_{Z,y}$ is a complete intersection of dimension
at least three.
\end{enumerate}
Here complete intersection means that its completion is a quotient
of a regular complete local ring by a regular sequence.
Restriction is an equivalence
\[
                    \operatorname{FEt}(Z)
                 \xrightarrow{\ \sim\ }\operatorname{FEt}(U).
\]
Consequently $\pi_1(U)\to\pi_1(Z)$ is an isomorphism, and these
groups are topologically finitely generated.
In particular, D.2(ii) holds if $Y$ has codimension at least three
and $Z$ is a local complete intersection at every point of $Y$.
Neither finiteness of $Y$ nor any condition on the singularities
retained in $U$ is required.
\end{proposition}
\begin{proof}
We use local purity in its precise essential-surjectivity form:
every finite \'etale cover of the punctured spectrum of a regular
local ring of dimension at least two, or of a complete-intersection
local ring of dimension at least three, extends to a finite \'etale
cover of its spectrum. These are, respectively,
\cite[Lemma 58.21.3(1), tag 0BMA; Proposition 58.25.3, tag 0BPD]{StacksDLocalPurity}.

The hypotheses imply that $Y$ contains no generic or codimension-one
point of $Z$, so $U$ is dense. Since $Z$ is normal and connected,
it is integral. Let $V\to U$ be a connected finite \'etale cover.
Normalize $Z$ in the finite separable extension $k(V)/k(Z)$, and
write $q:\overline V\to Z$ for the result. The normalization is
finite because $Z$ is of finite type over a field. Its restriction
to $U$ is $V$.

Let $B\subset Z$ be the image of the non-\'etale locus of $q$.
It is closed because $q$ is finite, and it is contained in $Y$.
Suppose $B\ne\varnothing$ and choose a generic point $b$ of an
irreducible component of $B$. Put $A=\mathcal O_{Z,b}$.
Every proper generization of $b$ lies outside $B$. Thus the base
change of $q$ to the punctured spectrum of $A$ is finite \'etale.
By one of the two local-purity statements it extends to a finite
\'etale $A$-scheme $W$.

The ring $A$ is normal, and $W$ is normal. Every component of $W$
dominates $\operatorname{Spec}A$: a finite \'etale morphism is open
and closed on each component, and the local base is connected.
Both $W$ and $\overline V\times_Z\operatorname{Spec}A$ are therefore
the integral closure of $A$ in the same finite product of generic
fields. They are isomorphic. It follows that $q$ is \'etale at
every point above $b$, contrary to $b\in B$. Hence $B$ is empty
and $V$ extends. Disconnected covers extend component by component.

Restriction is fully faithful as well: a morphism between two
finite \'etale covers of $Z$ is determined on their dense open
restrictions, and a morphism on those restrictions extends uniquely
by the integral-closure description in their generic fields.
This proves the asserted equivalence and the isomorphism on
fundamental groups, for any chosen geometric base point of $U$.
Finally, the fundamental group of a proper scheme over an
algebraically closed field is topologically finitely generated by
\cite[X, Theorem 2.9, pp.~273--274]{SGA1}. This proves the final
claim and the stated special case.
\end{proof}
